%% Conclusion Section
\section{Conclusion}
\label{sec:conclusion}

We have introduced the Adaptive Convergence Controller (ACC), a data-driven module
that replaces fixed training budgets and heuristic plateau counters in DRL-based
quantum state preparation with a principled, theoretically grounded stopping criterion.
ACC fits a saturating exponential model to the online training trajectory and identifies
three decision states — threshold met, flat unreachable, and converged — enabling
early termination of both successful and hopeless runs.

Integrated into QUASAR v14, ACC achieves fidelity $F \geq 0.99$ across all four
canonical entangled state families (GHZ, W, Cluster, Dicke-$k$) at 2 qubits with
30--98\% fewer training steps than fixed-budget baselines, validated across three
random seeds. The theoretical analysis in \cref{sec:theory} establishes that the
saturating exponential model arises from the Bellman contraction property of the
optimal value function near the target state, providing a first-principles
justification for ACC's convergence model.

The scaling conjecture $\tau(n) = \tau_0 \cdot \beta^{n-2}$ (\cref{eq:tau_scaling}),
if confirmed by the 3Q+ experiments in job 181423, would provide a predictive framework
for estimating the required training budget at any qubit count from measurements at
small qubit counts. This would be a significant step toward scalable DRL-based quantum
control on near-term hardware.

\subsection{Limitations and Future Work}

The current results are limited to 2 qubits, where the learning problem is relatively
tractable. The key open question is whether ACC's benefits persist at 3Q+ where the
Hilbert space dimension grows exponentially and the saturating exponential model may
break down. We are currently running experiments at 3Q and 4Q (job 181423) and will
update the results section upon completion.

A second limitation is that the theoretical analysis in \cref{sec:theory} relies on
local Lipschitz continuity and reward shaping assumptions that may not hold globally.
A more rigorous analysis using PAC-MDP theory \cite{kakade2003sample} or regret bounds
for entropy-regularised RL \cite{haarnoja2018soft} would strengthen the theoretical
contribution.

Finally, while ACC is presented here in the context of quantum state preparation, the
framework is algorithm-agnostic and applicable to any DRL problem with a saturating
reward signal. Future work will benchmark ACC against fixed-budget and ASHA baselines
on standard continuous control benchmarks (MuJoCo) and discrete action spaces (Atari),
establishing ACC as a general-purpose convergence controller for DRL.
